\subsection{BÀI TẬP TRẮC NGHIỆM}
\ind{PHẦN I.} \inden{Câu trắc nghiệm nhiều phương án lựa chọn. Mỗi câu hỏi học sinh chỉ chọn một phương án.}\\
\setcounter{ex}{0}
\Opensolutionfile{ans}[ans/1D6-B4-1]
\begin{ex}%[2D2B5-1]
	Phương trình $2^{2x+1}=32$ có nghiệm là
	\choice
	{$x=\dfrac{5}{2}$}
	{\True $x=2$}
	{$x=\dfrac{3}{2}$}
	{$x=3$}
	\loigiai{
		Ta có $2^{2x+1}=32 \Leftrightarrow 2x+1=5 \Leftrightarrow x=2$.}
\end{ex}


\begin{ex}%[TT, Lê Xoay, Vĩnh Phúc, 2018, L3]%[2D2Y5-1]%[Vinh Vo, Dự án (12EX-8)]
	Tích tất cả các nghiệm của phương trình $ 2^{x^2 + x } = 4 $ bằng
	\choice
	{$ 2 $}
	{$ 3 $}
	{\True $ -2 $}
	{$ -1 $}
	\loigiai{
		Ta có $ 2^{x^2 + x } = 4 \Leftrightarrow x^2 + x - 2 =0 \Leftrightarrow x = -2 \lor x = 1$.
	}	
\end{ex}

\begin{ex}
	Tọa độ giao điểm của đồ thị hàm số $y=2^{- x} + 3$ và đường thẳng $y=11$ là
	\choice
	{\True $\left(- 3; 11\right)$}
	{$\left(4; 11\right)$}
	{$\left(- 4; 11\right)$}
	{$\left(3; 11\right)$.}
	\loigiai{
		Phương trình hoành độ giao điểm $11=2^{- x} + 3\Leftrightarrow 2^{- x}=8\Leftrightarrow x= - 3$}
\end{ex}

\begin{ex}%[Đề HK2-2018, Chuyên Phan Ngọc Hiển-Cà Mau]%[Phùng Hoàng Em, dự án EX9]%[2D2B5-2]
	Biết rằng phương trình $2^{x^2-4x+2}=2^{x-4}$ có hai nghiệm phân biệt là $x_1$, $x_2$. Tính giá trị biểu thức $S=x_1^4+x_2^4$.
	\choice
	{$S=17$}
	{$S=257$}
	{\True $S=97$}
	{$S=92$}
	\loigiai{
		$2^{x^2-4x+2}=2^{x-4} \Leftrightarrow x^2-4x+2=x-4 \Leftrightarrow x^2-5x+6=0 \Leftrightarrow \hoac{&x_1=2\\&x_2=3.}$\\
		Suy ra $S=x_1^4+x_2^4=2^4+3^4=97$.}
\end{ex}

\begin{ex}%[HK2, THPT Phạm Công Bình, Vĩnh Phúc, 2018]%[Hà Lê, 12EX-9]%[2D2Y5-2]
	Tìm nghiệm của phương trình $5^{2018x}=\sqrt{5}^{2018}$.
	\choice
	{$x=1-\log_5 2$}
	{$x=-\log_5 2$}
	{\True $x=\dfrac{1}{2}$}
	{$x=2$}
	\loigiai{
		Ta có 
		\begin{eqnarray*}
			5^{2018x}=\sqrt{5}^{2018}&   \Leftrightarrow & 5^{2018x}={5}^{1009}\\
			& \Leftrightarrow & 2018x =1009 \\
			& \Leftrightarrow & x=\dfrac{1}{2}.
		\end{eqnarray*}
	}
\end{ex}

\begin{ex}%[Đề KSCL Toán 12 lần 2 năm 2017 - 2018, Phan Chu Trinh, Đắk Lắc]%[2D2B5-2]%[Cao Thành Thái, dự án 12EX-7]
	Tìm nghiệm của phương trình $9^{\sqrt{x - 1}} = \mathrm{e}^{\ln 81}$.
	\choice
	{\True $x = 5$}
	{$x = 4$}
	{$x = 6$}
	{$x = 17$}
	\loigiai
	{
		Điều kiện của phương trình $x - 1 \geq 0 \Leftrightarrow x \geq 1$.
		\begin{align*}
			& 9^{\sqrt{x - 1}} = \mathrm{e}^{\ln 81} \\
			\Leftrightarrow \ & 9^{\sqrt{x - 1}} = 81 \\
			\Leftrightarrow \ & \sqrt{x - 1} = 2 \\
			\Leftrightarrow \ & x - 1 = 4 \\
			\Leftrightarrow \ & x = 5.
		\end{align*}
		Đối chiếu điều kiện $x \geq 1$ ta nhận $x = 5$.\\
		Vậy phương trình đã cho có nghiệm duy nhất $x = 5$.
	}
\end{ex}

\begin{ex}%[2D2Y5-1]
	Tìm tập nghiệm $S$ của phương trình $\left(\dfrac{2}{3}\right)^{4x}=\left(\dfrac{3}{2}\right)^{2x-6}$
	\choice
	{$ S=\left\{{-1}\right\}$}
	{\True $ S=\left\{ 1\right\}$}
	{$ S=\left\{{-3}\right\}$}
	{$ S=\left\{ 3\right\}$}
	\loigiai{
		Phương trình tương đương $\left(\dfrac{2}{3}\right)^{4x}=\left(\dfrac{2}{3}\right)^{6-2x}\Leftrightarrow 4x=6-2x\Leftrightarrow x=1.$}
\end{ex}

\begin{ex} Tìm tất cả giá trị của tham số $m$ để phương trình $5^x-1-m=0$ có nghiệm.
	\choice
	{ $m>0$}
	{\True $m>-1$}
	{$m<0$}
	{$m<-1$}
	\loigiai{Phương trình tương đương $5^x=m+1$.\\
		Do $5^x>0$ nên phương trình này có nghiệm khi $m+1>0 \Leftrightarrow m>-1$.}
\end{ex} 

\begin{ex}%[2D2B5-1]
	Nghiệm của phương trình $\log_2x=3$ là
	\choice
	{$9$}
	{$6$}
	{\True $8$}
	{$5$}
	\loigiai{$\log_2x=3\Leftrightarrow x=2^3=8$.
	}
\end{ex}

\begin{ex}%[GHK2 - Thoại Ngọc Hầu - 2018]%[2D2B5-1]%[Đỗ Đường Hiếu]
	Tìm nghiệm của phương trình $\log_{64}(x+1)=\dfrac{1}{2}$.
	\choice
	{$-1$}
	{$4$}
	{\True $7$}
	{$-\dfrac{1}{2}$}
	\loigiai{
		Ta có $\log_{64}(x+1)=\dfrac{1}{2}\Leftrightarrow x+1=64^{\frac{1}{2}}\Leftrightarrow x+1=8 \Leftrightarrow x=7$.
	}
\end{ex}

\begin{ex}%[Phan Minh Tâm, dự án EX11-2018]%[2-BGD-102 - Đề thi chính thức THPT QG môn Toán mã 102 năm 2018]%[2D2Y5-1]
	Tập nghiệm của phương trình $\log_2(x^2-1)=3$ là
	\choice
	{\True $\{-3;3\}$}
	{$\{-3\}$}
	{$\{3\}$}
	{$\{-\sqrt{10};\sqrt{10}\}$}
	\loigiai{
		Ta có $\log_2(x^2-1)=3 \Leftrightarrow x^2-1=2^3\Leftrightarrow \hoac{&x=3\\&x=-3}$.\\
		Vậy tập nghiệm của phương trình đã cho là $ \lbrace -3; 3\rbrace $.
	}
\end{ex}


\begin{ex}%[GHK2, Trần Hưng Đạo,Hà Nội, 2018]%[2D2B5-1]%[Đỗ Đường Hiếu, dự án (12EX-7)]
	Số giao điểm của đồ thị hàm số $y=\log_2(x^2+3x)$ và đường thẳng $y=2$ là
	\choice
	{$0$}
	{\True $2$}
	{$1$}
	{$3$}
	\loigiai{
		Phương trình hoành độ giao điểm:
		$$\log_2(x^2+3x)=2\Leftrightarrow x^2+3x=4 \Leftrightarrow \hoac{&x=1\\&x=-4}.$$
		Vậy, đồ thị hàm số $y=\log_2(x^2+3x)$ và đường thẳng $y=2$ cắt nhau tại $2$ điểm phân biệt.
	}
\end{ex}

\begin{ex}%[GHK2, THPT Nhã Nam - Bắc Giang, 2018]%[2D2B5-2]%[Dương BùiĐức, 12EX8]
	Số nghiệm của phương trình $\log_{3}(x^{2}-6)=\log_{3}(x-2)+1$ là
	\choice
	{\True $1$}
	{$3$}
	{$2$}
	{$0$}
	\loigiai{
		Điều kiện xác định $\heva{&x-2>0\\ &x^{2}-6>0}\Leftrightarrow x>\sqrt{6}$.\\
		Ta có $\log_{3}(x^{2}-6)=\log_{3}(x-2)+1\Leftrightarrow x^{2}-6=3(x-2)\Leftrightarrow\hoac{&x=0\ (\mbox{loại})\\ &x=3\ (\mbox{thỏa mãn}).}$
	}
\end{ex}


\begin{ex}%[HK2, THPT Nguyễn Trãi - Hà Nội, 2018]%[Dương BùiĐức, 12EX9]%[2D2B5-2]
	Giải phương trình $ \log_{4}(x+1)+\log_{4}(x-3)=3 $.
	\choice
	{$ x=1\pm2\sqrt{17} $}
	{\True $ x=1+2\sqrt{17}  $}
	{$ x=5  $}
	{$ x=33  $}
	\loigiai{
		Điều kiện xác định $ \heva{&x+1>0\\ &x-3>0}\Leftrightarrow x>3 $.\\
		Ta có $ \log_{4}(x+1)+\log_{4}(x-3)=3\Leftrightarrow (x+1)(x-3)=64\Leftrightarrow x^{2}-2x-67=0\Leftrightarrow x=1\pm 2\sqrt{17} $. Kết hợp với điều kiện xác định ta có nghiệm của phương trình là $ x=1+2\sqrt{17}  $.
	}
\end{ex}


\begin{ex}%[2D2Y6-1]
	Tìm tập nghiệm $S$ của bất phương trình $\left( \dfrac{1}{2}\right) ^x>8$.
	\choice
	{$S=\left( -3;+\infty\right) $}
	{$S=\left( -\infty;3\right) $}
	{\True $S=\left( -\infty;-3\right) $}
	{$S=\left( 3;+\infty\right) $}
	\loigiai{
		Vì cơ số $\dfrac{1}{2}<1$ nên $\left( \dfrac{1}{2}\right) ^x>8 \Leftrightarrow x< \log_{\frac{1}{2}}8 \Leftrightarrow x<-3$.	
	}
\end{ex}

\begin{ex}%[2D2Y6-1]
	Tập nghiệm của bất phương trình $2^{2x}< 2^{x + 6}$ là
	\choice
	{$(0;6)$}
	{\True $( - \infty ; 6)$}
	{$( 0 ; 64 )$}
	{$(6 ; + \infty)$}
	\loigiai{
		Ta có 
		$2^{2x}< 2^{x + 6}\Leftrightarrow 2x<x+6 \Leftrightarrow x < 6.$
	}
	
\end{ex}

\begin{ex}%[2D2B6]
	Bất phương trình $\left(\dfrac{1}{2}\right)^{x^2+4x}>\dfrac{1}{32}$ có tập nghiệm là $S=(a;b)$. Khi đó giá trị $b-a$ là
	\choice
	{$4$}
	{$2$}
	{\True $6$}
	{$8$}
	\loigiai{
		BPT $\Leftrightarrow x^2+4x < 5 \Leftrightarrow x \in (-5;1) \Rightarrow T = 6.$
	}
\end{ex}

\begin{ex}%[2D2B6]
	Số nghiệm nguyên của bất phương trình ${\left(\sqrt{2}\right)}^{x^2-2x}\leqslant {\left(\sqrt{2}\right)}^3$ là
	\choice
	{3}
	{2}
	{\True 5}
	{4}
	\loigiai{
		${\left(\sqrt{2}\right)}^{x^2-2x}\leqslant {\left(\sqrt{2}\right)}^3 \Leftrightarrow x^2-2x-3\leqslant 0 \Leftrightarrow -1\leqslant x\leqslant 3$. \\
		Do $x$ nguyên, suy ra $x\in \{-1;0;1;2;3\}$. Vậy bất phương trình có 5 nghiệm nguyên.}
\end{ex}

\begin{ex}%[2D2Y6-1]
	Tập nghiệm của bất phương trình $\log _{0,3}(3x-2)\geq 0$ là
	\choice
	{$\left(\dfrac{2}{3};+\infty\right)$}
	{$\left(\dfrac{2}{3}; 1\right)$}
	{\True $\left(\dfrac{2}{3};,1\right]$}
	{$(2;+\infty)$}
	\loigiai{
		Ta có:
		$\log _{0,3}(3x-2)\geq 0 \Leftrightarrow \heva{& 3x-2>0\\&3x-2\leq 1}\Leftrightarrow \heva{& x>\dfrac{2}{3}\\&x\leq 1}\Leftrightarrow \dfrac{2}{3}<x\leq 1$.\\
		Tập nghiệm của bất phương trình là $\left(\dfrac{2}{3}; 1\right]$.}
\end{ex}

\begin{ex}%[2D2B6]
	Tập nghiệm của bất phương trình $\log_{0{,}5}\left(x-3\right)<\log_{0{,}5}\left(x^2-4x+3\right)$ là
	\choice
	{$\left(3;+\infty \right)$}
	{$\mathbb{R}$}
	{\True $\varnothing $}
	{$\left(2;3\right)$}
	\loigiai
	{
		Điều kiện: $\heva{
			& x-3>0 \\ 
			& x^2-4x+3>0}\Leftrightarrow \heva{
			& x>3 \\ 
			& x\in \left(-\infty;1\right)\cup \left(3;+\infty \right)}\Leftrightarrow x>3$. \\
		Bpt $\Leftrightarrow x-3>x^2-4x+3\Leftrightarrow x^2-5x+6<0\Leftrightarrow 2<x<3$. \\
		Kết hợp với điều kiện ta có nghiệm của bất phương trình là $S=\varnothing$.
	}
\end{ex}


\begin{ex}%[2D2Y6-1]
	Tập nghiệm của bất phương trình $\log(2x-1)\le \log x$ là
	\choice
	{$\left[\dfrac{1}{2};1\right]$}
	{$\left(-\infty;1\right]$}
	{\True $\left(\dfrac{1}{2};1\right]$}
	{$(0;1]$}
	\loigiai{
		$\log(2x-1)\le \log x$ $\Leftrightarrow 0<2x-1\le x$ $\Leftrightarrow \dfrac{1}{2}<x\le 1$.\\
		Vậy tập nghiệm của bất phương trình là $\left(\dfrac{1}{2};1\right]$
	}
\end{ex}

\begin{ex}%[2D2B6-1]
	Tập nghiệm $S$ của bất phương trình $\log_{\frac{1}{2}} (x^2-5x+7)>0$ là
	\choice
	{$S=(-\infty;2)$}
	{\True $S=(2;3)$}
	{$S=(3;+\infty)$}
	{$S=(-\infty;2)\cup (3;+\infty)$}
	\loigiai
	{
		Ta có
		\[\log_{\frac{1}{2}} (x^2-5x+7)>0 \Leftrightarrow \left\{\begin{aligned}&x^2-5x+7>0 \\&x^2-5x+7<1\end{aligned}\right. \Leftrightarrow x^2-5x+6<0 \Leftrightarrow 2<x<3.\]
		Vậy tập nghiệm của bất phương trình đã cho là $S=(2;3)$.
	}
\end{ex}

\begin{ex} Tìm tất cả giá trị của tham số $m$ để bất phương trình $3^x+1 \ge m$ có tập nghiệm là $\mathbb{R}$.
	\choice
	{$m<0$}
	{\True $m \le 1$ }
	{$m \le 0$}
	{$m>1$}
	\loigiai{Bất phương trình tương đương $3^x \ge m-1$.\\
		Yêu cầu bài toán $\Leftrightarrow m-1 \le 0 \Leftrightarrow m \le 1$.}
\end{ex}


\begin{ex}%[Thi thử, Sở GDĐT Kon Tum, 2020]%[Trần Chiến, 12EX11]%[2D2B6-6]%
	Ông An gửi tiền vào ngân hàng với  thể thức lãi kép theo công thức $T_n=A(1+r)^n$, trong đó $A$ là số tiền gửi ban đầu, $r$ là lãi suất, $n$ là số kì hạn, $T_n$ là số tiền cả gốc lẫn lãi sau $n$ kì hạn gửi. Nếu ông An gửi $100$ triệu đồng vào ngân hàng với lãi suất $6{,}5\%$/năm và không rút tiền gốc lẫn lãi định kì thì sau bao nhiêu năm ông ấy nhận được số tiền ít nhất là $250$ triệu đồng.
	\choice
	{$3$}
	{$10$}
	{\True $15$}
	{$8$}
	\loigiai{
		Từ giả thiết ta có bất phương trình $100(1+6{,}5\%)^n\geq 250\Leftrightarrow n\geq 14{,}55$.\\
		Vậy sau $15$ năm, ông An nhận được số tiền ít nhất là $250$ triệu đồng.
	}
\end{ex}

\begin{ex}%[Đề tập huấn, Sở GD - ĐT tỉnh Quảng Bình, 2019]%[Nguyễn Tiến, dự án 12EX5]%[2D2T5-6]
	Một người gửi tiết kiệm với lãi suất $0{,}7\%$/tháng và lãi hàng tháng được nhập vào vốn, hỏi sau bao nhiêu tháng người đó thu được gấp đôi số tiền ban đầu?
	\choice
	{$96$}
	{$97$}
	{$99$}
	{\True $100$}
	\loigiai{
		Gọi $x$ là số tiền gửi ban đầu $(x>0)$.\\
		%		Do lãi suất $1$ năm là $8{,}4\%$ nên lãi suất của $1$ tháng là $0{,}7\%$.\\
		Số tiền sau tháng thứ nhất là $1{,}007x$.\\
		Số tiền sau tháng thứ hai là $(1{,}007)^2x$.\\
		Số tiền sau tháng thứ $n$ là $(1{,}007)^nx$.\\
		Theo giả thiết ta có $(1{,}007)^nx=2x\Leftrightarrow(1{,}007)^n=2\Leftrightarrow n=100.$
	}
\end{ex}

\Closesolutionfile{ans}

\ind{PHẦN II.} \inden{Câu trắc nghiệm đúng sai. Trong mỗi ý a), b), c), d) ở mỗi câu, học sinh chọn đúng hoặc sai.}\\
\Opensolutionfile{ans}[ans/1D6-B4-2]

\begin{ex}
	Cho hai hàm số $f(x)=2^{x^2-3x+4}$ và $g(x)=\left(\dfrac{1}{2}\right)^{x-10}$. Xét tính đúng - sai các khẳng định sau
	\choiceTF
	{\True Phương trình $f(x)=1$ vô nghiệm}
	{\True Phương trình $g(x)=1$ có nghiệm là số thực dương}
	{Phương trình $f(x)=g(x)$ có tập nghiệm là $S=\{-2;3\}$}
	{Bất phương trình $f(x)\leq g(x)$ có tập nghiệm là $S=\{1-\sqrt{7};1+\sqrt{7}\}$}
	\loigiai{
	}
\end{ex}

\begin{ex}%[1D6H4-2]
	Cho phương trình $\log (x-1)^2=\log (x+1)$. Xét tính đúng - sai các khẳng định sau
	\choiceTF
	{Điều kiện $x>1$}
	{Phương trình đã cho có chung tập nghiệm với phương trình $x^2-3x+\dfrac{9}{4}=0$}
	{\True Tổng các nghiệm của phương trình bằng $3$}
	{\True Biết phương trình có hai nghiệm $x_1,x_2,(x_1<x_2 )$. Khi đó $3$ số $x_1;x_2;6$ tạo thành một cấp số cộng}
	\loigiai{
		Điều kiện: $\heva{
			& (x-1)^2>0 \\
			& x+1>0}\Leftrightarrow -1<x \ne 1$.\\
		$\log (x-1)^2=\log (x+1) \Rightarrow(x-1)^2=x+1 \Leftrightarrow x^2-3x=0 \Leftrightarrow \hoac{
			&x=0 \\
			&x=3.}$\\
		So điều kiện phương trình có tập nghiệm là $S=\{0;3\}$.
		\begin{itemchoice}
			\itemch Sai. Do điều kiện $-1<x\ne 1$.
			\itemch Sai. Do $x=3$ không là nghiệm phương trình $x^2-3x+\dfrac{9}{4}=0$.
			\itemch Đúng. Do Tổng hai nghiệm là $0+3=3$.
			\itemch Đúng. Do $0,3,6$ lập thành cấp số cộng có công sai bằng $3$.
		\end{itemchoice}
	}
\end{ex}

\begin{ex}%[1D6H4-5]%[Tổ Latex - Strong]%[Tổ 18 - Đợt 13 - CK2 - Lớp 11 - KNTT]%[Thái Thịnh]
	Cho hàm số $f(x)=\dfrac{\log_{3} (x+25) -4}{5^{x^2}-125^x}$. Xét tính đúng - sai các khẳng định sau
	\choiceTF
	{Hàm số $f(x)$ có tập xác định là $\mathscr{D}=\left[-25; +\infty\right)\setminus \{0;3\}$}
	{\True Phương trình $f(x)=0$ có duy nhất một nghiệm}
	{\True Tập nghiệm của bất phương trình $f(x)\ge 0$ là $(0;3)\cup \left[56; +\infty\right)$}
	{\True Tập nghiệm của bất phương trình $f(x)<0$ có tất cả $52$ số nguyên dương}
	\loigiai{
		\begin{itemchoice}
			\itemch Xét mệnh đề: Hàm số $f(x)$ có tập xác định là $\mathscr{D}=\left[-25; +\infty\right)\setminus \{0;3\}$.\\
			Điều kiện xác định $\heva{&x+25>0\\&5^{x^2}-125^x\ne 0}
			\Leftrightarrow\heva{&x>-25\\&5^{x^2} \ne 5^{3x}}  
			\Leftrightarrow\heva{&x>-25\\&x\ne0\\&x\ne3.}$\\ 
			Nên tập xác định của hàm số đã cho là $\mathscr{D}=\left(-25; +\infty\right)\setminus \{0;3\}$.\\
			Vậy mệnh đề này sai.
			\itemch Xét mệnh đề: Phương trình $f(x)=0$ có duy nhất một nghiệm.\\
			Ta có $f(x)=0 \Leftrightarrow \log_{3} (x+25)-4=0 \Leftrightarrow \log_{3} (x+25)=4 \Leftrightarrow x=56$ (nhận).\\
			Vậy mệnh đề này đúng.
			\itemch Xét mệnh đề: Tập nghiệm của bất phương trình $f(x)\ge 0$ là $(0;3)\cup \left[56; +\infty\right)$.\\
			Ta có $f(x)\ge 0 \Leftrightarrow \dfrac{\log_{3} (x+25) -4}{5^{x^2}-125^x}\ge 0$.\\
			Điều kiện xác định $\heva{&x+25>0\\&5^{x^2}-125^x\ne 0}
			\Leftrightarrow\heva{&x>-25\\&5^{x^2} \ne 5^{3x}}  
			\Leftrightarrow\heva{&x>-25\\&x\ne0\\&x\ne3.}$\\ 
			Ta có $\log_{3} (x+25)-4=0 \Leftrightarrow \log_{3} (x+25)=4 \Leftrightarrow x=56$ (nhận).\\
			Sử dụng phương pháp xét dấu khoảng, ta có bảng xét dấu vế trái của bất phương trình đã cho
			\begin{center}
				\begin{tikzpicture}
					\tkzTabSetup[doubledistance= 2pt]
					\tkzTabInit[lgt=2,espcl=2] % tùy chọn
					{$x$/0.8, Vế trái/0.8} % cột 1
					{$-25$,$0$,$3$,$56$,$+\infty$} % hàng 1 cột 2
					\tkzTabLine{ , - , d , + , d , - , 0,+, } % hàng 2 cột 2
				\end{tikzpicture}
			\end{center}
			Suy ra $f(x)\ge 0 \Leftrightarrow x\in(0;3)\cup \left[56; +\infty\right)$. Vậy mệnh đề này đúng.
			\itemch Xét mệnh đề: Tập nghiệm của bất phương trình $f(x)<0$ có tất cả $52$ số nguyên dương.\\
			Ta có $f(x)< 0 \Leftrightarrow \dfrac{\log_{3} (x+25) -4}{5^{x^2}-125^x}< 0$.\\
			Điều kiện xác định $\heva{&x+25>0\\&5^{x^2}-125^x\ne 0}
			\Leftrightarrow\heva{&x>-25\\&5^{x^2} \ne 5^{3x}}  
			\Leftrightarrow\heva{&x>-25\\&x\ne0\\&x\ne3.}$\\ 
			Ta có $\log_{3} (x+25)-4=0 \Leftrightarrow \log_{3} (x+25)=4 \Leftrightarrow x=56 $ (nhận).\\
			Sử dụng phương pháp xét dấu khoảng, ta có bảng xét dấu vế trái của bất phương trình đã cho:
			\begin{center}
				\begin{tikzpicture}
					\tkzTabSetup[doubledistance= 2pt]
					\tkzTabInit[lgt=2,espcl=2] % tùy chọn
					{$x$/0.8, Vế trái/0.8} % cột 1
					{$-25$,$0$,$3$,$56$,$+\infty$} % hàng 1 cột 2
					\tkzTabLine{ , - , d , + , d , - , 0,+, } % hàng 2 cột 2
				\end{tikzpicture}
			\end{center}
			Suy ra $f(x)<0 \Leftrightarrow x\in(-25;0)\cup (3;56)$.\\
			Nên tập nghiệm của bất phương trình $f(x)<0$ có tất cả $52$ số nguyên dương.\\
			Vậy mệnh đề này đúng.
		\end{itemchoice}
	}
\end{ex}

\Closesolutionfile{ans}

\ind{PHẦN III.} \inden{Câu trắc nghiệm trả lời ngắn.}\\
\Opensolutionfile{ans}[ans/1D6-B4-3]

\begin{ex}
	Cho hai hàm số $f(x)=x(x-1)(x-2)\cdots(x-10)$ và $g(x)=2(x-1)(x-2)\cdots(x-10)$. Tính tổng tất cả các nghiệm của phương trình $2^{f(x)}=2^{g(x)}$.
	\\
	\shortans[oly]{$55$}
	\loigiai{
		$2^{f(x)}=2^{g(x)} \Leftrightarrow x(x-1)(x-2)\cdots(x-10)=2(x-1)(x-2)\cdots(x-10) \Leftrightarrow  x \in \{1;2;\cdots;10\}$.\\
		Tổng các nghiệm là $S=1+2+...+10=55$.
	}
\end{ex}

\begin{ex}
Sự tăng trưởng của một loài vi khuẩn được tính theo công thức $f(t)=A \cdot e^n$, trong đó $A$ là số lượng vi khuẩn ban đầu, $r$ là tỷ lệ tăng trưởng ( $r>0$ ), $t$ (tính theo giờ) là thời gian tăng trưởng. Biết số vi khuẩn ban đầu có 1000 con và sau 10 giờ là 5000 con. Hỏi sau ít nhât bao nhiêu giờ thì số lượng vi khuẩn tăng gấp 10 lần? (Làm tròn kêt quá đên hàng phần chục).\\
\shortans[oly]{$14,3$}
	\loigiai{
	Theo giả thiết thì
	$$
		5000=1000 e^{10 r} \Leftrightarrow e^{10 r}=5 \Leftrightarrow r=\frac{\ln 5}{10} .
		$$
	Gọi t (giờ) là thời gian cần tìm để sô lượng vi khuẩn tăng gấp 10 lần Do đó: $10000=1000 e^{r t} \Leftrightarrow e^{r t}=10 \Leftrightarrow r t=\ln 10 \Leftrightarrow t=\frac{\ln 10}{r}=\frac{\ln 10}{\frac{\ln 5}{10}}=10 \log _s 10 \approx 14,3$
	Vậy sau khoảng 14,3 giờ thì số lượng vi khuẩn sẽ tăng gãp 10 lần.
	}
\end{ex}


\begin{ex}%[1D6V4-6]%Câu20
	Một nghiên cứu cho thấy một nhóm học sinh được cho xem cùng danh sách các loài động vật và được kiểm tra lại xem họ nhớ bao nhiêu \% mỗi tháng. Sau t tháng, khả năng nhớ trung bình của nhóm học sinh được cho bới công thức $M(t)=75-20\ln (t+1 ),t\ge 0$ (đơn vị \%). Hỏi sau khoảng bao nhiêu tháng thì nhóm học sinh đó nhớ được danh sách dưới 10\%?\\
	\shortans[oly]{$25$}
	\loigiai{
		Theo công thức tính tỉ lệ phần trăm ta cần tìm $t$ thỏa mãn:
		$$75-20\ln (t+1 )<10 \Leftrightarrow \ln (t+1 )>3{,}25\Leftrightarrow t+1>25{,}79\Leftrightarrow t>24{,}79.$$ }
\end{ex}

\begin{ex}%[2D2B6-6]%
	Để quảng bá cho sản phẩm $A$, một công ty dự định tổ chức quảng cáo theo hình thức quảng cáo trên truyền hình. Nghiên cứu của công ty cho thấy: Nếu sau $n$ quảng cáo được phát thì tỉ lệ người xem quảng cáo đó mua sản phẩm $A$ tuân theo công thức: $P(n)=\dfrac{1}{1+49 \cdot \mathrm{e}^{-0,015n}}$. Hỏi cần ít nhất bao nhiêu lần quảng cáo để tỉ lệ người xem mua sản phẩm đạt trên $30\%$?
	\\
	\shortans[oly]{$203$}
	\loigiai{
		Để số người mua sản phẩm đạt trên $30\%$
		\allowdisplaybreaks
		\begin{eqnarray*}
			\Rightarrow P(n)>0,3 &\Leftrightarrow& \dfrac{1}{1+49 \cdot \mathrm{e}^{-0,015n}}>0,3\\
			&\Leftrightarrow& 1+49 \cdot \mathrm{e}^{-0,015n}<\dfrac{1}{0,3}\\
			&\Leftrightarrow& \mathrm{e}^{-0,015n}<\dfrac{1}{21}\\
			&\Leftrightarrow& \mathrm{e}^{0,015n}>21\\
			&\Leftrightarrow& 0,015n>\ln 21 \approx 202,97.
		\end{eqnarray*}
		Vậy phải có ít nhất $203$ lần quảng cáo.
	}
\end{ex}

\begin{ex}%[2D2B6-6]
	Anh Nam muốn mua một ngôi nhà trị giá $ 500 $ triệu đồng sau $ 3 $ năm nữa. Biết rằng lãi suất hàng năm vẫn không đổi là $ 8\% $ một năm. Vậy ngay từ bây giờ số tiền ít nhất anh Nam phải gửi tiết kiệm vào ngân hàng theo thể thức lãi kép để có đủ tiền mua nhà (kết quả làm tròn đến hàng triệu) là bao nhiêu?\\
	\shortans[oly]{$397$}
	\loigiai{
		Gọi $ x $ là số tiền anh Nam phải giải tiết kiệm, đơn vị triệu đồng. Khi đó sau 3 năm, theo hình thức lãi kép thì số tiền anh Nam nhận được là $ x\left(1+ \dfrac{8}{100} \right)^3 $. \
		Vậy ta cần tìm $ x $ là số nguyên nhỏ nhất thỏa mãn điều kiện
		\[ x\left(1+ \dfrac{8}{100} \right)^3 \ge 500 \Leftrightarrow x \ge \dfrac{500}{\left(1+ \dfrac{8}{100} \right)^3} \approx 396,9. \]
		Vậy $ x = 397 $ triệu đồng.	
	}
\end{ex}


\begin{ex}%[Đề KSCL T12 lần 3, Triệu Quang Phục-Hưng Yên 2018]%[2D2K6]
	Một người sử dụng xe máy có giá trị ban đầu là $40$ triệu đồng. Sau mỗi năm, giá trị xe giảm $10\%$ so với năm trước đó. Hỏi sau ít nhất bao nhiêu năm thì giá trị xe nhỏ hơn $12$ triệu đồng?\\
	\shortans[oly]{$12$}
	\loigiai{
		Giá trị xe sau năm thứ nhất là:\\
		$40-40\cdot \dfrac{1}{10}$ triệu đồng.\\
		Giá trị xe sau năm thứ hai là:\\
		$40-40\cdot \dfrac{1}{10}-\left(40-40\cdot\dfrac{1}{10}\right)\dfrac{1}{10}=\left(40-40\cdot\dfrac{1}{10}\right)\left(1-\dfrac{1}{10}\right)$ triệu đồng.\\
		Giá trị xe sau năm thứ ba là:\\
		$\left(40-40\cdot\dfrac{1}{10}\right)\left(1-\dfrac{1}{10}\right)-\left(40-40\cdot\dfrac{1}{10}\right)\left(1-\dfrac{1}{10}\right)\dfrac{1}{10}=\left(40-40\cdot\dfrac{1}{10}\right)\left(1-\dfrac{1}{10}\right)^2$\\
		$\ldots$\\
		Suy ra giá xe sau năm thứ $n$ là:\\
		$\left(40-40\cdot\dfrac{1}{10}\right)\left(1-\dfrac{1}{10}\right)^{n-1}$.\\
		$\left(40-40\cdot\dfrac{1}{10}\right)\left(1-\dfrac{1}{10}\right)^{n-1}<12\Leftrightarrow 36\cdot\left(\dfrac{9}{10}\right)^{n-1}<12\Leftrightarrow \left(\dfrac{9}{10}\right)^{n-1}<\dfrac{1}{3}\Leftrightarrow n-1>\log_{\frac{9}{10}}\dfrac{1}{3}$ \\
		$\Leftrightarrow n>1-\log_{\frac{9}{10}}3>11{,}4$.\\
		Vậy sau ít nhất $12$ năm thì giá xe nhỏ hơn $12$ triệu đồng.
	}
\end{ex}

\begin{ex}%[ThanhTa, dự án EX-THPTQG-20-103]%[2D2K6-6]
	(QG.2020 lần 1 -- Mã đề 103).  Trong năm $2019$, diện tích rừng trồng mới của tỉnh $A$ là $900$ ha. Giả sử diện tích rừng trồng mới của tỉnh $A$ mỗi năm tiếp theo đều tăng $6\%$ so với diện tích rừng trồng mới của năm liền trước. Kể từ sau năm $2019$, năm nào là năm đầu tiên tỉnh $A$ có diện tích rừng trồng mới trong năm đó đạt trên $1700$ ha? \\
	\shortans[oly]{$2030$}
	\loigiai{
		Yêu cầu bài toán $\Leftrightarrow 900\cdot (1+0{,}06)^n >1700 \Leftrightarrow n>\log_{1{,}06}\dfrac{1700}{900}\approx 10{,}9$.\\
		Vậy sau $11$ năm tức năm $2030$ thì tỉnh $A$ có diện tích rừng trồng mới trong năm đó đạt trên $1700$ ha.
	}
\end{ex}

\Closesolutionfile{ans}

